% !TeX root = ../main.tex

\section{Conclusion}

Windowed CUSUM provides a locally adaptive framework for reconstructing recurrent stand disturbances from long Sentinel-1 time series. Track-specific trajectories preserve acquisition-geometry consistency, overlapping windows repeatedly renew the temporal reference, and window-specific permutation calibration evaluates each CUSUM departure relative to the variability of the same local context. Under spatial out-of-fold evaluation, Windowed CUSUM achieved an F1 of 0.797, compared with 0.680 for CCDC VV and 0.675 for decrease-filtered ReCuSum. The advantage arose primarily from higher recurrent-event recall. Component analysis further identified local-reference renewal as the main source of sensitivity, while central event admission and local statistical calibration recovered precision.

Applied to 3,498 validated balsa polygons, the locked detector reconstructed 6,658 disturbances and 3,509 consecutive inter-disturbance intervals during 2017--2025, with a median recurrence of 36 months. More broadly, the study shows that geometry-consistent SAR trajectories combined with local temporal and statistical adaptation can extend disturbance monitoring from isolated change detection toward reconstruction of recurrent plantation disturbance histories in cloud-prone regions.